\section*{Evaluasi}
\addcontentsline{toc}{section}{Evaluasi}

\begin{enumerate}[label=\arabic*.,leftmargin=*]
  \item Definisikan tata bahasa ambigu dan berikan satu contoh.
  \item Jelaskan dua pohon urai pada dangling else.
  \item Tuliskan aturan else terdekat.
  \item Hilangkan left recursion dari \(A \to A a \mid A b \mid c\).
  \item Hilangkan left recursion tidak langsung dari
        \(A \to B a \mid c\), \(B \to A b \mid d\).
  \item Lakukan left factoring pada
        \(S \to \code{if}\ E\ \code{then}\ S\ \code{else}\ S
        \mid \code{if}\ E\ \code{then}\ S\).
  \item Jelaskan mengapa parser top-down tidak dapat menangani left
        recursion.
  \item Diberikan grammar yang masih memiliki konflik, lakukan debug
        sampai bebas konflik dan buktikan kesetaraan bahasa dengan lima
        string uji.
\end{enumerate}

\subsection*{Kunci dan Rubrik Singkat}
\begin{itemize}[leftmargin=*]
  \item Soal 1--3: ketepatan konsep (70\%), contoh (30\%).
  \item Soal 4--6: kebenaran transformasi (80\%), kelengkapan (20\%).
  \item Soal 7: penjelasan (70\%), argumentasi (30\%).
  \item Soal 8: hasil debug (60\%), bukti kesetaraan (40\%).
\end{itemize}
